\begin{frame}{$\lambda$ 를 바꾸면 부호까지 바뀐다}
  \small
  \begin{tabular}{@{}lcll@{}}
    \toprule
    $\lambda$ & 가중계수 $(1,\gamma\lambda,\gamma^2\lambda^2)$
      & $A_0^{\mathrm{GAE}}$ & 읽기 \\
    \midrule
    0.0 & $(1,\ 0,\ 0)$ & $0.26$ & $\delta_0$ 하나만 — ``1스텝만
      보고'' (TD(0)) \\
    0.5 & $(1,\ 0.45,\ 0.2025)$ &
      $0.26+0.45(0.14)+0.2025(-0.60)=+0.20$ & 중간 — 세 항 다
      포함 \\
    0.95 & $(1,\ 0.855,\ 0.731)$ & $-0.059$ & 거의 MC \\
    1.0 & $(1,\ 0.9,\ 0.81)$ &
      $G_0-V(s_0)=1.9-2.0=-0.10$ & 전체 리턴 기준 (MC) \\
    \bottomrule
  \end{tabular}
  \vskip0.8em
  \begin{itemize}
    \item \textbf{같은 데이터}($\delta_0,\delta_1,\delta_2$)인데
      $\lambda$ 만 바꾸면 $+0.26$(좋았다)에서 $-0.10$(나빴다)까지
      \kb{부호까지} 바뀜.
    \item $\lambda$ 가 ``\kb{신뢰하는 지평선의 길이}''이라는
      뜻: 짧은 지평선($\lambda\approx0$)은 ``지금 당장의 TD
      오차만 믿는다''(분산 적음, $V$ 의 오차가 편향으로),
      긴 지평선($\lambda=1$)은 ``에피소드 전체의 실제 리턴을
      믿는다''(편향 없음, 리턴의 흔들림이 분산으로).
    \item 실전 기본값 $\lambda=0.95$ = ``긴 지평선에 거의
      의존하되, 가장 먼 꼬리만 5\% 정도로 얕은 신뢰''의 절충.
  \end{itemize}
\end{frame}
